\documentclass[11pt]{article}
\usepackage[T1]{fontenc}
\usepackage[utf8]{inputenc}
\usepackage{lmodern}
\usepackage[margin=1in]{geometry}
\usepackage{amsmath,amssymb,amsthm,mathtools}
\usepackage{microtype,booktabs,longtable,array}
\usepackage[dvipsnames]{xcolor}
\usepackage{enumitem,fancyhdr,titlesec}
\usepackage{listings,needspace}
\usepackage[most]{tcolorbox}
\usepackage[colorlinks=true,linkcolor=MidnightBlue,citecolor=MidnightBlue,urlcolor=teal]{hyperref}
\definecolor{ink}{HTML}{19354C}
\definecolor{accent}{HTML}{147D92}
\definecolor{pale}{HTML}{F0F6F8}
\titleformat{\section}{\Large\bfseries\color{ink}}{\thesection}{0.7em}{}
\titleformat{\subsection}{\large\bfseries\color{ink}}{\thesubsection}{0.7em}{}
\pagestyle{fancy}
\fancyhf{}
\fancyhead[L]{\small\color{ink}Polynomial divisibility of factorial ratios}
\fancyhead[R]{\small Research article}
\fancyfoot[C]{\thepage}
\setlength{\headheight}{14pt}
\setlength{\parskip}{4pt}
\setlength{\parindent}{1.3em}
\setlength{\emergencystretch}{2em}
\allowdisplaybreaks[2]
\newtheorem{theorem}{Theorem}[section]
\newtheorem{lemma}[theorem]{Lemma}
\newtheorem{proposition}[theorem]{Proposition}
\newtheorem{corollary}[theorem]{Corollary}
\theoremstyle{definition}
\newtheorem{definition}[theorem]{Definition}
\newtheorem{example}[theorem]{Example}
\theoremstyle{remark}
\newtheorem{remark}[theorem]{Remark}
\newcommand{\Z}{\mathbb Z}
\newcommand{\Q}{\mathbb Q}
\newcommand{\N}{\mathbb N}
\newcommand{\vp}{v_p}
\newcommand{\floor}[1]{\left\lfloor #1\right\rfloor}
\newcommand{\ind}{\mathbf 1}
\newcommand{\lcm}{\operatorname{lcm}}
\newcommand{\sgn}{\operatorname{sgn}}
\newcommand{\dd}{\,\mathrm d}
\newcommand{\A}{A}
\newcommand{\B}{B}
\newcommand{\LL}{\mathcal L}
\newcommand{\hyp}[3]{{}_{8}F_{7}\!\left(\begin{matrix}#1\\#2\end{matrix};#3\right)}
\lstset{basicstyle=\ttfamily\small,breaklines=true,columns=fullflexible,
  frame=single,rulecolor=\color{accent},backgroundcolor=\color{pale},
  showstringspaces=false}
\tcbset{colback=pale,colframe=accent,boxrule=0.6pt,arc=1mm,left=3mm,right=3mm}
\hypersetup{pdftitle={Polynomial divisibility of height-one factorial ratios},
 pdfauthor={OpenAI ChatGPT},pdfsubject={Bala conjectures, OEIS A211417 and A295431}}

\begin{document}
\begin{titlepage}
\vspace*{0.5cm}
{\color{accent}\large RESEARCH ARTICLE\par}
\vspace{0.6cm}
{\color{ink}\Huge\bfseries Polynomial divisibility\par of height-one\par factorial ratios\par}
\vspace{0.6cm}
{\Large Proofs of Bala's uniform-multiplier conjectures,\par
 a complete rational-root criterion,\par and certified optimal constants\par}
\vspace{0.8cm}
{\large A research attempt with reproducible exact certificates\par}
{\large September 19, 2026\par}
\vspace{0.8cm}
\begin{tcolorbox}[title=Main outcome,fonttitle=\bfseries]
For
\[
 A(n)=\frac{(30n)!\,n!}{(15n)!\,(10n)!\,(6n)!},
\]
we prove the all-$r$ divisibility conjectures still recorded under OEIS
A211417. We also classify every polynomial splitting over $\Q$ that can
uniformly divide $A(n)$ after multiplication by one fixed integer.
The analogous conjectures for OEIS A295431 follow from the same theorem.
\end{tcolorbox}
\vspace{0.4cm}
\noindent\textbf{Abstract.}
Let $R(n)$ be a balanced integral factorial ratio and let $P\in\Z[X]$
split over $\Q$. We characterize the existence of a nonzero integer $C$
with $CR(n)/P(n)\in\Z$ for every nonnegative integer $n$ at which the
quotient is defined. The criterion compares the multiplicity of each
rational root of $P$ with one-sided values of the Landau step function.
We prove both directions and give an explicit multiplier. At height one,
the criterion reduces to the signs of cyclotomic exponents and excludes
all repeated roots. Applications settle uniform-product conjectures of
Peter Bala, including explicit all-$r$ bounds. Finite weighted digit
graphs determine the eight particular proposed constants and prove each
optimal. The associated integer sequences admit explicit recurrences,
hypergeometric generating functions, and all-order asymptotic expansions.
The archive contains exact certificates, a separate verifier, regression
tests, and an admissible-denominator atlas for all 52 sporadic examples.
\vfill
\end{titlepage}

\begingroup
\small
\setlength{\parskip}{0pt}
\tableofcontents
\endgroup
\newpage

\section{The selected problem and its status}
\subsection{Two exceptional integer sequences}
Consider the factorial ratios
\begin{equation}\label{eq:AB}
 A(n)=\frac{(30n)!\,n!}{(15n)!\,(10n)!\,(6n)!},
 \qquad
 B(n)=\frac{(12n)!\,n!}{(6n)!\,(4n)!\,(3n)!},\qquad n\geq0.
\end{equation}
They are OEIS A211417 and A295431, respectively. Both belong to the
52 sporadic height-one integral factorial ratios in the classification
of Bober, following the work of Vasyunin and the hypergeometric approach
of Rodriguez-Villegas and Beukers--Heckman \cite{oeisA,oeisB,bober}.
Our arguments require no classification of hypergeometric monodromy.

Bala's August 2025 comments on A211417 ask for a fixed multiplier,
depending on $k$ and $r$ but not on $n$, that clears the denominators in
\begin{equation}\label{eq:balaPositive}
 \frac{A(n)}{\displaystyle\prod_{\substack{1\leq i\leq r\\(i,k)=1}}(kn+i)},
 \qquad k\in\{2,3,5\}.
\end{equation}
There is also the negative-offset family
\begin{equation}\label{eq:balaNegative}
 \frac{A(n)}{\displaystyle\prod_{\substack{1\leq i\leq r\\(i,30)=1}}(30n-i)}.
\end{equation}
The current entry still labels these all-$r$ assertions conjectural
\cite{oeisA}. Its one-factor case $A(n)/(30n-1)$ is different: the entry
records a proof in June 2026, and Adamczewski's August 2026 report also
lists it as proved \cite{oeisA,adamczewski}. We use it only as a previously
known special case recovered by a broader argument.

The companion A295431 entry poses similar uniform-product questions for
$B(n)$, using positive slopes $1,2,3$ and negative factors $12n-i$ with
$(i,12)=1$ \cite{oeisB}. The selected problem is therefore not merely to
reprove integrality of \eqref{eq:AB}, or to reproduce the solved first
negative-offset case. It is to explain and prove the uniform-product
phenomenon.

\subsection{Results obtained}
Write
\[
 \LL_s=\lcm(1,2,\ldots,s),\qquad \LL_1=1,
 \qquad m_k(r)=\#\{1\leq i\leq r:(i,k)=1\}.
\]
Our first concrete conclusions are
\begin{align}
 &\LL_{30r/k}^{\,m_k(r)}
   \frac{A(n)}{\prod_{1\leq i\leq r,\,(i,k)=1}(kn+i)}\in\Z
   &&(k=2,3,5),\label{eq:explicitC}\\
 &\LL_r^{\,m_{30}(r)}
   \frac{A(n)}{\prod_{1\leq i\leq r,\,(i,30)=1}(30n-i)}\in\Z.
   \label{eq:explicitD}
\end{align}
These hold for every $r\geq1$ and $n\geq0$. The indices $30r/k$ are
integers. The displayed constants are deliberately explicit rather than
claimed to be smallest. The corresponding negative-offset formula for
$B$ uses $\LL_r^{m_{12}(r)}$.

The structural result is stronger. A polynomial splitting over $\Q$
can divide $A(n)$ uniformly up to one fixed integer multiplier precisely
when it has distinct, nonzero rational roots, and the reduced
denominator of each negative root lies in $\{1,2,3,5\}$, while the reduced
denominator of each positive root is $30$. For $B$, these sets become
$\{1,2,3\}$ and $\{12\}$. This also proves that the positive-offset
coprimality restrictions can be removed, at the cost of changing the
constant.

The eight particular constants in the two OEIS entries are correct and
optimal:
\begin{center}
\begin{tabular}{lll}
\toprule
Ratio & Dividing polynomial & Least positive multiplier\\
\midrule
$A$ & $2n+1$ & $7$\\
$A$ & $3n+1$ & $1$\\
$A$ & $5n+1$ & $1$\\
$A$ & $(2n+1)(3n+1)(5n+1)$ & $42$\\
$B$ & $n+1$ & $385$\\
$B$ & $2n+1$ & $5$\\
$B$ & $3n+1$ & $1$\\
$B$ & $(n+1)(2n+1)(3n+1)$ & $770$\\
\bottomrule
\end{tabular}
\end{center}
Existence for these eight cases follows already from the structural
proof. Their exact constants are established by the finite certificates
in Section~\ref{sec:sharp}.

\subsection{What is and is not asserted about novelty}
The source check was made on September 19, 2026. It distinguishes the
still-conjectural all-$r$ statements from the explicitly solved
$30n-1$ case. The machinery of factorial valuations and Landau step
functions is classical. We give a complete proof of the additional
polynomial-divisibility statements here, but do not claim that an
exhaustive literature search has ruled out every equivalent formulation
of the general criterion. In particular, this manuscript's mathematical
claims are separable from a claim of publication priority.

\section{Factorial valuations and the Landau function}
\subsection{Definitions and a self-contained integrality criterion}
Let $a_1,\ldots,a_u$ and $b_1,\ldots,b_v$ be positive integers satisfying
\begin{equation}\label{eq:balance}
 \sum_{i=1}^u a_i=\sum_{j=1}^v b_j.
\end{equation}
Set
\begin{equation}\label{eq:RDelta}
 R(n)=\frac{\prod_{i=1}^u(a_i n)!}{\prod_{j=1}^v(b_j n)!},\qquad
 \Delta(x)=\sum_{i=1}^u\floor{a_i x}-\sum_{j=1}^v\floor{b_j x},
 \qquad h=v-u.
\end{equation}
The integer $h$ is the height. The balance condition makes $\Delta$
periodic with period one, since the terms linear in an integer
translation cancel. It is a right-continuous integer-valued step
function. Put $M=\max(a_i,b_j)$.

For a prime $p$, $\vp(q)$ denotes its exponent in a nonzero rational
number $q$; it may be negative. Counting the multiples of $p,p^2,\ldots$
in a factorial gives
\begin{equation}\label{eq:legendre}
 \vp(N!)=\sum_{e\geq1}\floor{N/p^e},
 \qquad
 \vp(R(n))=\sum_{e\geq1}\Delta(n/p^e).
\end{equation}
All these sums are finite. A nonzero rational number is an integer if
and only if all its prime valuations are nonnegative.

\begin{lemma}[One-variable Landau criterion]\label{lem:landau}
Under \eqref{eq:balance}, $R(n)\in\Z$ for every $n\geq0$ if and only if
$\Delta(x)\geq0$ for every real $x$.
\end{lemma}
\begin{proof}
Nonnegativity gives integrality immediately from \eqref{eq:legendre}.
Conversely, suppose $\Delta$ is negative somewhere. Right continuity
and the finite set of discontinuities in a period give a nonempty open
interval $(\alpha,\beta)\subset(0,1)$ on which it is negative. Choose
an arbitrarily large prime $p$ and an integer $n$ with
$\alpha<n/p<\beta$; such an $n$ exists once $p(\beta-\alpha)>1$.
For $p>M$ we have $0\leq n/p^e<1/M$ for $e\geq2$. On this last
interval every floor in \eqref{eq:RDelta} is zero. Thus
$\vp(R(n))=\Delta(n/p)<0$, contradicting integrality. Periodicity
handles every real $x$.
\end{proof}
This is the relevant specialization of Landau's theorem; Bober gives
its general setting and consequences \cite[Section 3]{bober}. Bala's
note supplies a direct proof for $A$ and its floor-factorial extension
\cite{bala-note}.

\subsection{Reflection and jumps}
For $k\geq1$, define
\begin{equation}\label{eq:J}
 J(k)=\#\{i:k\mid a_i\}-\#\{j:k\mid b_j\}.
\end{equation}
If $(t,k)=1$, the jump at $t/k$ is
\begin{equation}\label{eq:jump}
 \Delta(t/k)-\Delta((t/k)^-)=J(k).
\end{equation}
Indeed, $\floor{c x}$ increases at $t/k$ precisely when $k\mid c$.
Here $\Delta(x^-)$ is the value just to the left of $x$; the value just
to the right is $\Delta(x)$ itself.

Away from the discontinuities,
\begin{equation}\label{eq:reflection}
 \Delta(x)+\Delta(1-x)=h.
\end{equation}
To see this, use $\floor{c x}+\floor{c(1-x)}=c-1$ when $cx$ is not
an integer and sum over the parameter lists. At a reduced rational
point the endpoint version is
\begin{align}
 \Delta(t/k)+\Delta(1-t/k)&=h+J(k),\label{eq:refRight}\\
 \Delta((t/k)^-)+\Delta((1-t/k)^-)&=h-J(k).\label{eq:refLeft}
\end{align}
These follow either by counting the integral $cx$ separately or by
subtracting the two jumps. The formulas remain valid at integers by
periodicity.

\begin{corollary}\label{cor:01}
For an integral balanced ratio of height one, every value and one-sided
value of $\Delta$ belongs to $\{0,1\}$, and $J(k)\in\{-1,0,1\}$.
\end{corollary}
\begin{proof}
At a point avoiding the discontinuities, the two nonnegative integers
in \eqref{eq:reflection} have sum one. Right continuity gives the same
range at every point, and left limits have the same range. Equation
\eqref{eq:jump} then bounds the jumps.
\end{proof}

\section{A complete criterion for polynomials with rational roots}
\subsection{The one-sided capacity}
Every nonzero polynomial in $\Z[X]$ that splits over $\Q$ can be written
\begin{equation}\label{eq:Pfactor}
 P(X)=c\prod_{j=1}^s(k_jX+\ell_j)^{m_j},
\end{equation}
where $c\in\Z\setminus\{0\}$, $k_j>0$, $(k_j,\ell_j)=1$, $m_j\geq1$,
and the roots $-\ell_j/k_j$ are pairwise distinct. This is primitive
factorization over $\Z$ after collecting equal roots. Write
$d=\sum_jm_j$. A constant polynomial is allowed by taking $s=0$.

For $\ell\ne0$, define the capacity
\begin{equation}\label{eq:capacity}
 \kappa_R(k,\ell)=
 \begin{cases}
 \displaystyle\min_{\substack{1\leq t\leq k\\(t,k)=1}}
       \Delta((t/k)^-),&\ell>0,\\[2ex]
 \displaystyle\min_{\substack{1\leq t\leq k\\(t,k)=1}}
       \Delta(t/k),&\ell<0.
 \end{cases}
\end{equation}
For $k=1$, the only representative is $t=1$; thus the left capacity is
$h$ and the right capacity is zero. We also set the capacity of the
primitive zero-root factor $X$ equal to zero.

The sign matters. A positive offset $kn+\ell$ forces $n/p^e$ to
approach a reduced rational point from the left. A negative offset
forces approach from the right. The size of $|\ell|$ affects the
exceptional small prime powers, but not the capacity.

\begin{theorem}[Polynomial divisibility criterion]\label{thm:criterion}
Assume $R$ is balanced and integral. For $P$ as in \eqref{eq:Pfactor},
the following are equivalent:
\begin{enumerate}[label=(\roman*),itemsep=2pt]
\item There is a positive integer $C$ such that
$CR(n)/P(n)\in\Z$ for every $n\geq0$ with $P(n)\ne0$.
\item The polynomial has no zero root and
$m_j\leq\kappa_R(k_j,\ell_j)$ for every $j$.
\end{enumerate}
When (ii) holds and $s>0$, one valid choice is
\begin{equation}\label{eq:Cbound}
 C=|c|\LL_T^{\,d},\qquad
 T=\max\left(1,\ M\max_j|\ell_j|,
    \max_{i<j}|k_i\ell_j-k_j\ell_i|\right),
\end{equation}
with the last maximum omitted if $s=1$.
For $s=0$, $C=|c|$ works.
\end{theorem}

\subsection{Proof of sufficiency}
We isolate the large-prime-power part of the proof.

\begin{lemma}\label{lem:large}
Suppose $\ell\ne0$, $(k,\ell)=1$, and $q=p^e>M|\ell|$. If
$q\mid kn+\ell$ and $kn+\ell\ne0$, then
\[
 \Delta(n/q)\geq\kappa_R(k,\ell).
\]
\end{lemma}
\begin{proof}
If $p\mid k$, then $p\nmid\ell$, so divisibility is impossible.
Otherwise put $kn+\ell=tq$. Since $q>|\ell|$ and $n\geq0$, a nonzero
multiple $tq=kn+\ell$ cannot be negative; hence $t\geq1$.
Reducing modulo $k$ gives $(t,k)=1$, since both $q$ and $\ell$ are
coprime to $k$. We have
\begin{equation}\label{eq:approach}
 \frac{n}{q}=\frac{t}{k}-\frac{\ell}{kq}.
\end{equation}
Every discontinuity of $\Delta$ has the form $v/a_i$ or $v/b_j$.
A discontinuity different from $t/k$ is at distance at least
$1/(kM)$: a nonzero difference of these rational numbers has numerator
at least one after placing it over denominator $ka_i$ or $kb_j$.
The displacement in \eqref{eq:approach} is strictly smaller than this.
Thus $\Delta(n/q)$ equals the left value at $t/k$ if $\ell>0$, and
the right value if $\ell<0$. Periodicity reduces $t$ to the
representatives in \eqref{eq:capacity}.
\end{proof}

\begin{proof}[Sufficiency in Theorem~\ref{thm:criterion}]
Let $Q(X)=\prod_j(k_jX+\ell_j)^{m_j}$. For a prime power $q$, put
\[
 E_q(n)=\sum_{j=1}^s m_j\ind_{q\mid k_jn+\ell_j}.
\]
If $q>T$, it cannot divide two distinct primitive factors. Otherwise
it would divide the nonzero integer
$k_i\ell_j-k_j\ell_i$, whose absolute value is at most $T$.
Consequently, for $q>T$, either $E_q(n)=0$ or exactly one root
contributes, with its multiplicity $m_j$. Lemma~\ref{lem:large} and
(ii) give
\begin{equation}\label{eq:termwiseHigh}
 \Delta(n/q)-E_q(n)\geq0\qquad(q>T).
\end{equation}
For $q\leq T$, Landau nonnegativity and $E_q(n)\leq d$ give the lower
bound $-d$. Therefore, for every prime $p$,
\begin{align*}
 \vp(R(n)/Q(n))
 &=\sum_{e\geq1}\bigl(\Delta(n/p^e)-E_{p^e}(n)\bigr)\\
 &\geq-d\#\{e\geq1:p^e\leq T\}
 =-d\vp(\LL_T).
\end{align*}
This proves $\LL_T^dR(n)/Q(n)\in\Z$. Since $P=cQ$, multiplying
by $|c|$ only adds the harmless sign $|c|/c$. No unproved assertion
about an infinite computational range enters the argument.
\end{proof}

\subsection{Proof of necessity}
The converse explains why the allowed slopes are restrictive.

\begin{proof}[Necessity in Theorem~\ref{thm:criterion}]
First suppose $P$ has a zero root. Primitivity makes its factor $X$.
For any sufficiently large prime $p$, take $n=p$. Then
$\Delta(n/p)=\Delta(1)=0$, and $\Delta(n/p^e)=0$ for $e\geq2$.
Hence $\vp(R(p))=0$, while $p\mid P(p)$, except for finitely many
$p$ that could make $P(p)=0$. No fixed multiplier can supply every
such prime.

Now fix a nonzero root with primitive factor $kX+\ell$, multiplicity
$m$, and capacity $\kappa<m$. Choose a reduced representative $t/k$
at which the relevant one-sided value is $\kappa$. Dirichlet's theorem
provides infinitely many primes in the progression
\begin{equation}\label{eq:dirichletclass}
 p\equiv \ell\,t^{-1}\pmod k;
\end{equation}
this is an eligible class because $(t,k)=(\ell,k)=1$. For $k=1$, it
means all primes. We use only the infinitude assertion of Dirichlet's
theorem \cite[Theorem 4.2]{kedlaya}.

For such a sufficiently large prime set
\begin{equation}\label{eq:witnessn}
 n=\frac{tp-\ell}{k}.
\end{equation}
Then $n$ is a nonnegative integer and $kn+\ell=tp$. Taking $p>k$
ensures $p\nmid t$, so the factor has valuation exactly one, and its
$m$th power contributes $m$ to $\vp(P(n))$. As $p$ tends to infinity,
$n/p=t/k-\ell/(kp)$ approaches from the required side; thus
$\Delta(n/p)=\kappa$ eventually. For all sufficiently large $p$,
$Mn<p^2$, so every term with $e\geq2$ vanishes. It follows that
\[
 \vp(R(n))=\kappa<m\leq\vp(P(n)).
\]
The values \eqref{eq:witnessn} tend to infinity, so finitely many zeros
of $P$ do not obstruct the construction. Infinitely many different
primes appear in reduced denominators of $R(n)/P(n)$, ruling out every
fixed nonzero multiplier.
\end{proof}

\subsection{Consequences and computational meaning}
The criterion is effective. All relevant values of $\Delta$ and all
jumps can be computed with exact integer or rational arithmetic. It
separates the problem into two very different tasks: deciding whether
\emph{some} multiplier exists, and finding the \emph{least} multiplier.
The first is answered by a finite root-by-root test. The second must
account for compensation between valuation contributions at different
powers of the same small prime.

The multiplicity condition is also essential. The fact that a sequence
is divisible by one linear expression, up to a fixed constant, does
not imply the same for its square. Higher-height ratios can support
larger multiplicities: for example $A(n)^2$ supports $(2n+1)^2$, but
not $(2n+1)^3$. This follows either from the theorem or from doubling
the Landau function.

\section{The height-one simplification}
\begin{theorem}[Jump-sign classification]\label{thm:heightone}
Suppose $R$ is balanced, integral, and of height one. For a primitive
factor $kX+\ell$ with $\ell\ne0$, its capacity is one exactly when
\begin{equation}\label{eq:signcriterion}
 J(k)=-1\quad\hbox{if }\ell>0,
 \qquad
 J(k)=+1\quad\hbox{if }\ell<0.
\end{equation}
Otherwise its capacity is zero. Thus the admissible polynomials are
exactly those with distinct nonzero roots whose primitive factors
satisfy \eqref{eq:signcriterion}.
\end{theorem}
\begin{proof}
By Corollary~\ref{cor:01}, every value is zero or one. If $J(k)=-1$,
then \eqref{eq:jump} forces the left value to be one and the right value
to be zero at \emph{every} reduced $t/k$. Conversely, if all left
values are one, pair $t/k$ with $1-t/k$ in \eqref{eq:refLeft}. Their
sum is two, so $1-J(k)=2$ and $J(k)=-1$.
The same argument with \eqref{eq:refRight} proves that all right
values are one exactly when $J(k)=+1$. This includes the self-paired
points and $k=1$, so no endpoint exception is needed. All capacities
are at most one, and Theorem~\ref{thm:criterion} finishes the proof.
\end{proof}

\subsection{Cyclotomic interpretation}
Let $\Phi_k(z)$ be the $k$th cyclotomic polynomial. Factorization of
$z^m-1$ gives the rational-function identity
\begin{equation}\label{eq:cyclo}
 \frac{\prod_i(z^{a_i}-1)}{\prod_j(z^{b_j}-1)}
 =\prod_{k\geq1}\Phi_k(z)^{J(k)}.
\end{equation}
Only finitely many exponents are nonzero. At height one the allowed
positive-offset denominators are exactly the cyclotomic factors
remaining downstairs after cancellation in \eqref{eq:cyclo}; the
allowed negative-offset denominators are those remaining upstairs.
This is an interpretation of the proved criterion, not an additional
monodromy assumption.

For the binomial example $R(n)=\binom{2n}{n}$, the sets are
$\{1\}$ and $\{2\}$. Accordingly, fixed products of distinct
$n+i$ can be divided out up to constants, as can fixed products of
distinct $2n-i$ with $i$ odd, but a repeated such root is impossible.
For the 52 sporadic ratios, Appendix~\ref{app:atlas} computes the two
sets explicitly from their parameter lists.

\section{Resolution of the uniform-product conjectures}
\subsection{The two exact slope lists}
For $A$, direct cancellation in \eqref{eq:cyclo} gives
\begin{equation}\label{eq:cycloA}
 \frac{(z^{30}-1)(z-1)}{(z^{15}-1)(z^{10}-1)(z^6-1)}
 =\frac{\Phi_{30}(z)}{\Phi_1(z)\Phi_2(z)\Phi_3(z)\Phi_5(z)}.
\end{equation}
For $B$,
\begin{equation}\label{eq:cycloB}
 \frac{(z^{12}-1)(z-1)}{(z^6-1)(z^4-1)(z^3-1)}
 =\frac{\Phi_{12}(z)}{\Phi_1(z)\Phi_2(z)\Phi_3(z)}.
\end{equation}
These identities prove the slope lists announced in Section 1, provided
integrality is known. For completeness, it can be checked particularly
simply here. On the grid interval $j/30\leq x<(j+1)/30$,
\[
 \Delta_A(x)=j-\floor{j/2}-\floor{j/3}-\floor{j/5},\quad 0\leq j<30.
\]
The three subtracted floors have sum at most
$\floor{31j/30}=j$. This proves nonnegativity and hence integrality
of $A$. For $B$ the corresponding formula is
\[
 \Delta_B(x)=j-\floor{j/2}-\floor{j/3}-\floor{j/4},\quad 0\leq j<12,
\]
and the sum is at most $\floor{13j/12}=j$. Height one then gives the
upper bound one in each case.

\begin{corollary}\label{cor:classificationAB}
Let $P\in\Z[X]\setminus\{0\}$ split over $\Q$. There exists a fixed
positive integer $C$ with $CA(n)/P(n)\in\Z$ wherever defined on
$\Z_{\geq0}$ if and only if $P$ has distinct nonzero roots, every
negative root in lowest terms has denominator $1,2,3$, or $5$, and
every positive root in lowest terms has denominator $30$.
For $B$, replace these lists by $1,2,3$ and $12$.
\end{corollary}

\subsection{Sharper explicit bounds for the prescribed products}
The general bound \eqref{eq:Cbound} already establishes existence.
The grid structure gives the cleaner formulas
\eqref{eq:explicitC}--\eqref{eq:explicitD}.

\begin{proposition}\label{prop:positive}
For $k\in\{2,3,5\}$, let $I=\{1\leq i\leq r:(i,k)=1\}$ and
$m=|I|$. Then the multiplier $\LL_{30r/k}^{m}$ works for
$A(n)/\prod_{i\in I}(kn+i)$.
\end{proposition}
\begin{proof}
If $q=p^e>30r/k$ divides $kn+i$, then $p\nmid k$ and
$n/q=t/k-i/(kq)$ with $(t,k)=1$. The left grid interval of width
$1/30$ at $t/k$ has $\Delta_A=1$, since $J_A(k)=-1$. The displacement
is smaller than $1/30$, so $\Delta_A(n/q)=1$.
Moreover, $q>r$ cannot divide two factors, since their difference is
a nonzero integer $i-j$ of absolute value less than $r$.
For smaller $q$ the total deficit is at worst $m$. The same valuation
sum used in Theorem~\ref{thm:criterion} gives the proposed multiplier.
\end{proof}

\begin{proposition}\label{prop:negative}
Let $(R,M)$ be either $(A,30)$ or $(B,12)$ and let
$I=\{1\leq i\leq r:(i,M)=1\}$. Then
\begin{equation}\label{eq:negativeboth}
 \LL_r^{|I|}\frac{R(n)}{\prod_{i\in I}(Mn-i)}\in\Z
 \qquad(n\geq0).
\end{equation}
\end{proposition}
\begin{proof}
No denominator factor is zero, since $i$ is coprime to $M>1$.
If $q>r$ divides $Mn-i$, then $p\nmid M$ and
$n/q=t/M+i/(Mq)$ with $(t,M)=1$. The right grid interval of width
$1/M$ at $t/M$ has value one, because $J_R(M)=+1$. Hence the
corresponding Landau term pays for this divisor. Two factors cannot
both be divisible by $q>r$, because their difference has absolute
value less than $r$. The low prime powers contribute a deficit of at
most $|I|$ each, exactly as before.
\end{proof}
For $r=1$ the multiplier is one. Thus the proposition recovers the
previously proved $A(n)/(30n-1)$ theorem without depending on its proof.
The purpose of the proposition is its uniform extension to every $r$.
Negative values at $n=0$ cause no problem: divisibility and valuations
use the absolute value of a nonzero denominator.

\subsection{Removing a restriction on positive offsets}
For $R=A$ and $k\in\{1,2,3,5\}$, or for $R=B$ and
$k\in\{1,2,3\}$, consider the unrestricted product
\[
 P_{k,r}(n)=\prod_{i=1}^r(kn+i).
\]
Put $g_i=(k,i)$ and reduce its factors to
$(k/g_i)n+i/g_i$. Their roots are still distinct, and their reduced
slopes divide $k$. Every such divisor is in the corresponding
positive-offset list. Consequently a uniform multiplier exists.
A fully explicit choice, with $M=30$ or $12$ respectively, is
\begin{equation}\label{eq:unrestricted}
 \left(\prod_{i=1}^r g_i\right)\LL_{Mr}^{r}.
\end{equation}
Indeed, the primitive offsets have absolute value at most $r$, and
their pairwise resultants are
$k(j-i)/(g_i g_j)$, of absolute value less than $kr\leq Mr$.
Theorem~\ref{thm:criterion} applies, with the product of the $g_i$
accounting for the content of the polynomial.

In contrast, the negative-offset coprimality restrictions are
structural. Reducing $30n-i$ when $(30,i)>1$ gives a proper divisor of
$30$ as its slope, none of which is in the negative-offset list for
$A$. A uniform multiplier for that one factor is impossible. The same
statement holds for $12n-i$ and $B$.

\subsection{An explicit infinite obstruction family}
The slope $4$ is not allowed for positive offsets of $A$. This failure
can be exhibited without relying on a table of numerical experiments.
Let $p>30$ be a prime with $p\equiv3\pmod4$, and set
\[
 n=\frac{3p-1}{4}.
\]
Then $4n+1=3p$, whereas
$n/p=3/4-1/(4p)$ lies in $[22/30,23/30)$, where $\Delta_A=0$.
Also $30n<p^2$ for $p>30$, so all higher valuation terms vanish.
Thus $\vp(A(n))=0$ while $p\mid4n+1$. Infinitely many distinct such
primes exist, so no constant clears all denominators of $A(n)/(4n+1)$.
The first example in this family is $p=31$, $n=23$.

\section{Exact optimal multipliers by finite certificates}\label{sec:sharp}
\subsection{Why checking a large interval is not enough}
The large-prime-power proof gives a constant, but not necessarily the
smallest one. It is not valid to maximize the deficit at each exponent
independently and claim that the resulting sum is sharp: residues
at $p,p^2,p^3,\ldots$ are coupled by the base-$p$ digits of $n$.
For instance, an unfavorable contribution at $p$ can force a favorable
contribution at $p^2$ or later. This is why some quotients above need
no multiplier even though an individual Landau summand can fail to
cover the denominator at a small prime power.

We encode all possible digit continuations in a finite directed graph.
Its integer edge weights add to the exact valuation, not to a floating
point approximation. A feasible potential on the graph proves a lower
bound for every possible $n$ simultaneously.

\subsection{The state and transition rules}
For this section let $R=A$ or $B$, let $L=30$ or $12$ respectively,
and let
\[
 Q(n)=\prod_{k\in K}(kn+1),
\]
where $K$ is one of the eight slope sets in the result table. Write
$\delta(j)=\Delta_R(j/L)$ for $0\leq j<L$.
Process a base-$p$ expansion of $n$ from least significant digit to
most significant digit. After $e$ digits, let
$r_e=n\bmod p^e$ and set
\[
 j_e=\floor{Lr_e/p^e}.
\]
Because every factorial parameter divides $L$, the valuation summand
$\Delta_R(n/p^e)$ equals $\delta(j_e)$.

A factor $kn+1$ is called alive after $e$ digits when $p^e\mid kr_e+1$.
For an alive factor retain the integer carry
\[
 t_{k,e}=(kr_e+1)/p^e.
\]
Its possible values are $1,\ldots,k$. Once a factor is not alive,
it can never be alive again, since divisibility by $p^{e+1}$ implies
divisibility by $p^e$. Mark it by $\bot$ and forget its carry.
The initial state is
\[
 s_0=(0,1,\ldots,1).
\]
If the next digit is $d\in\{0,\ldots,p-1\}$, the next state satisfies
\begin{equation}\label{eq:digitJ}
 j'=\floor{(j+Ld)/p},
\end{equation}
and for each factor
\begin{equation}\label{eq:digitCarry}
 t'_k=
 \begin{cases}
 (t_k+kd)/p,&t_k\ne\bot\ \hbox{and }p\mid t_k+kd,\\
 \bot,&\hbox{otherwise}.
 \end{cases}
\end{equation}
To verify \eqref{eq:digitJ}, write $Lr_e/p^e=j+\varepsilon$ with
$0\leq\varepsilon<1$; dividing $j+Ld+\varepsilon$ by the integer $p$
does not alter the floor obtained from $j+Ld$.
The carry formula follows by replacing $r_e$ with $r_e+d p^e$.

Give this transition the weight
\begin{equation}\label{eq:edgeweight}
 w(s,d)=\delta(j')-\#\{k\in K:t'_k\ne\bot\}.
\end{equation}
Reading all the digits of $n$ and then enough zero digits reaches
\[
 s_*=(0,\bot,\ldots,\bot).
\]
Indeed $j_e=0$ when $p^e>Ln$, and no positive integer $kn+1$ is
divisible by every power of $p$. The terminal state remains fixed under
further zero digits and has zero edge weight. Equations
\eqref{eq:legendre} and \eqref{eq:edgeweight} show that the sum along
this finite path is exactly
\begin{equation}\label{eq:pathvalue}
 \vp(R(n)/Q(n)).
\end{equation}

\subsection{Potential certificates}
\begin{proposition}[Finite certificate rule]\label{prop:potential}
Let $S$ be a finite collection of states containing $s_0$ and $s_*$ and
closed under every digit transition. Suppose an integer function
$H:S\to\Z$ satisfies
\begin{equation}\label{eq:potential}
 H(s_*)=0,\qquad H(s)\leq w(s,d)+H(s')
 \quad\hbox{for every }s\in S\hbox{ and every digit }d.
\end{equation}
Then $\vp(R(n)/Q(n))\geq H(s_0)$ for every $n\geq0$.
If some explicit $n_0$ has valuation $H(s_0)$, the lower bound is sharp.
\end{proposition}
\begin{proof}
Apply the edge inequality successively along the path for $n$ described
above. The intermediate values of $H$ telescope, leaving
$H(s_0)\leq\sum w+H(s_*)=\vp(R(n)/Q(n))$. An equality witness proves
that no larger lower bound is possible.
\end{proof}
The generator enumerates reachable states and computes shortest costs
to the terminal by integer relaxation. This search is \emph{not} trusted
as the proof: a separate program reads the resulting state list and
potential and checks every closure condition and inequality in
\eqref{eq:potential}. It also checks the equality witness directly from
Legendre's formula. All prime certificates needed for a particular
quotient are checked for complete coverage.

\subsection{Results of the exact verification}
For $A$ the general cutoff \eqref{eq:Cbound} is $T=30$ for these
positive unit offsets; for $B$ it is $T=12$. Thus primes larger than
these bounds need no certificate: each of their valuation summands is
nonnegative after subtracting the denominator indicator. Below the
cutoff the exact minima are as follows. A zero beyond a row's cutoff
is both a lower bound from the theorem and an attained minimum at $n=0$.

\begin{center}
\setlength{\tabcolsep}{4pt}
\small
\begin{tabular}{lrrrrrrrrrrr}
\toprule
$(R;K)$ & $2$&$3$&$5$&$7$&$11$&$13$&$17$&$19$&$23$&$29$&$C_{\min}$\\
\midrule
$(A;2)$       &0&0&0&-1&0&0&0&0&0&0&7\\
$(A;3)$       &0&0&0&0&0&0&0&0&0&0&1\\
$(A;5)$       &0&0&0&0&0&0&0&0&0&0&1\\
$(A;2,3,5)$   &-1&-1&0&-1&0&0&0&0&0&0&42\\
$(B;1)$       &0&0&-1&-1&-1&0&0&0&0&0&385\\
$(B;2)$       &0&0&-1&0&0&0&0&0&0&0&5\\
$(B;3)$       &0&0&0&0&0&0&0&0&0&0&1\\
$(B;1,2,3)$   &-1&0&-1&-1&-1&0&0&0&0&0&770\\
\bottomrule
\end{tabular}
\end{center}
The headings $2,3,\ldots$ are primes, and the entries are
$\min_{n\geq0}\vp(R(n)/Q(n))$. The table is not inferred from testing
an interval: it is the output of the checked finite potentials.

The nonzero lower-bound requirements are all necessary, as witnessed
by the following small indices:
\begin{center}
\begin{tabular}{lll}
\toprule
$(R;K)$ & Prime & Index with valuation $-1$\\
\midrule
$(A;2)$ & $7$ & $10$\\
$(A;2,3,5)$ & $2,3,7$ & $3,16,10$, respectively\\
$(B;1)$ & $5,7,11$ & $9,20,10$, respectively\\
$(B;2)$ & $5$ & $2$\\
$(B;1,2,3)$ & $2,5,7,11$ & $1,2,20,10$, respectively\\
\bottomrule
\end{tabular}
\end{center}
For each prime, any valid multiplier must contain that prime to at
least the indicated deficit. Multiplying these independently necessary
prime powers gives the least multiplier; the indices need not be the
same for different primes. This proves optimality of all eight constants.

The delivered certificates comprise \textbf{60 prime graphs, 1,685
states, and 19,277 checked transition inequalities}. Their compact JSON
representation contains every state and every potential value, together
with equality witnesses. The separate verifier recomputes all
transitions instead of trusting stored edges.

\section{Integer sequences, recurrences, and generating functions}
\subsection{An integer sequence singled out by the optimal constant}
Define
\begin{equation}\label{eq:U}
 U_n=\frac{42A(n)}{(2n+1)(3n+1)(5n+1)}.
\end{equation}
The preceding section proves $U_n\in\Z_{>0}$, with the factor $42$
minimal for integrality at every $n\geq0$. Its first terms are
\begin{align*}
 U_0&=42,\\
 U_1&=45287852610,\\
 U_2&=5873158887743776788432,\\
 U_3&=1645528285131615053352729526247883.
\end{align*}
The archive supplies exact terms through $n=50$. We make no assertion
that this sequence lacks an existing OEIS entry; the contribution here
is its proved uniform integrality and optimal normalization.

For an arbitrary polynomial divisor $P$ covered by our theorem,
$u_n=CR(n)/P(n)$ has the explicit first-order recurrence
\begin{equation}\label{eq:genericrecurrence}
 \frac{u_{n+1}}{u_n}
 =\frac{P(n)}{P(n+1)}
 \frac{\displaystyle\prod_{i=1}^u\prod_{t=1}^{a_i}(a_i n+t)}
      {\displaystyle\prod_{j=1}^v\prod_{t=1}^{b_j}(b_j n+t)}.
\end{equation}
Cross-multiplication gives a recurrence with polynomial coefficients.
This follows directly by canceling consecutive factorials and is
valid wherever the quotient is defined. In our admissible cases
$P$ has no nonnegative integer root, so there is no missing index.

\subsection{A hypergeometric representation}
Set
\[
 \lambda=2^{14}3^9 5^5,
 \qquad
 \alpha=\frac1{30}(1,7,11,13,17,19,23,29),
\]
and let the seven denominator parameters be
\[
 \beta=\left(\frac12,\frac13,\frac23,
                    \frac15,\frac25,\frac35,\frac45\right).
\]
Cancellation of the linear factors in $A(n+1)/A(n)$ gives
\begin{equation}\label{eq:Apock}
 A(n)=\lambda^n\frac{\prod_{a\in\alpha}(a)_n}
                         {n!\prod_{b\in\beta}(b)_n},
\end{equation}
where $(a)_n=a(a+1)\cdots(a+n-1)$ and $(a)_0=1$.
One can verify the identity without invoking the gamma multiplication
formula: both sides have value one at $n=0$ and the same ratio of
successive terms. It is also consistent with the hypergeometric
description in A211417 \cite{oeisA}.

The elementary identity
\[
 \frac{1}{kn+1}=\frac{(1/k)_n}{(1+1/k)_n}
\]
shows that dividing by the three factors in \eqref{eq:U} shifts three
denominator parameters. Define
\[
 \beta'=\left(\frac32,\frac43,\frac23,
                      \frac65,\frac25,\frac35,\frac45\right).
\]
Then
\begin{equation}\label{eq:Upock}
 U_n=42\lambda^n\frac{\prod_{a\in\alpha}(a)_n}
                           {n!\prod_{b\in\beta'}(b)_n},
\end{equation}
and hence
\begin{equation}\label{eq:Uhypergeom}
 G(z):=\sum_{n\geq0}U_nz^n
 =42\hyp{\frac1{30},\frac7{30},\frac{11}{30},\frac{13}{30},
             \frac{17}{30},\frac{19}{30},\frac{23}{30},\frac{29}{30}}
 {\frac32,\frac43,\frac23,\frac65,\frac25,\frac35,\frac45}{\lambda z}.
\end{equation}
This is an equality of formal power series, and of analytic functions
for $|\lambda z|<1$. The coefficients also satisfy the convenient
recurrence
\begin{equation}\label{eq:Urec}
 (n+1)\prod_{b\in\beta'}(n+b)\,U_{n+1}
 =\lambda\prod_{a\in\alpha}(n+a)\,U_n,\qquad U_0=42.
\end{equation}
No recurrence is being used to define an unproved integer sequence:
integrality has already been established independently.

\subsection{Differential and integral transforms}
Let $\theta=z\,\mathrm d/\mathrm dz$ and
$F(z)=\sum_{n\geq0}A(n)z^n$. Coefficient extraction gives
\begin{equation}\label{eq:Euleroperator}
 (2\theta+1)(3\theta+1)(5\theta+1)G(z)=42F(z).
\end{equation}
Equivalently, within the common disk of convergence,
\begin{equation}\label{eq:integral}
 G(z)=42\int_0^1\!\int_0^1\!\int_0^1
       F(zx^2y^3t^5)\dd x\dd y\dd t.
\end{equation}
Termwise integration is justified by absolute convergence, and each
monomial contributes the three reciprocal factors in \eqref{eq:U}.

The first-order recurrence \eqref{eq:Urec} gives the differential equation
\begin{equation}\label{eq:hypode}
 \left[\theta\prod_{b\in\beta'}(\theta+b-1)
       -\lambda z\prod_{a\in\alpha}(\theta+a)\right]G(z)=0.
\end{equation}
For $n\geq1$, its $z^n$ coefficient is just the recurrence with index
$n-1$; its constant coefficient vanishes due to the initial factor
$\theta$. Thus $G$ is D-finite. More generally, polynomial division
of coefficients gives $P(\theta)\sum u_nz^n=C\sum R(n)z^n$.
These identities make the new integer divisibility statements useful
for recurrence and generating-function investigations, without needing
to guess a closed algebraic equation.

\section{All-order asymptotic expansions}
\subsection{The general factorial-ratio expansion}
The usual Stirling expansion for $\log\Gamma$, including remainder
control on the positive real axis, is recorded in DLMF Section 5.11
\cite{dlmf}. Applying it separately to the finitely many factorials
in $R(n)$ yields the following explicit form. Let
\[
 \Lambda=\frac{\prod_i a_i^{a_i}}{\prod_j b_j^{b_j}},\qquad
 K_R=(2\pi)^{-h/2}
       \sqrt{\frac{\prod_i a_i}{\prod_j b_j}},
\]
and, with $\mathsf B_{2q}$ the Bernoulli numbers, set
\begin{equation}\label{eq:sigma}
 \sigma_{2q-1}=
 \frac{\mathsf B_{2q}}{2q(2q-1)}
 \left(\sum_i a_i^{1-2q}-\sum_j b_j^{1-2q}\right).
\end{equation}
Then, for any fixed $N\geq1$,
\begin{equation}\label{eq:logasym}
 \log R(n)=n\log\Lambda-\frac h2\log n+\log K_R
 +\sum_{q=1}^{N}\frac{\sigma_{2q-1}}{n^{2q-1}}
 +O(n^{-2N-1}).
\end{equation}
The terms of size $n\log n$ and $n$ cancel by balance, the $\log n$
terms give $-h/2$, and the remaining terms are exactly those displayed.
Thus \eqref{eq:logasym} is a Poincare expansion to any prescribed order,
not a claim that its infinite series converges.

Suppose $P$ has degree $d$, positive leading coefficient $\rho$, and
factorization \eqref{eq:Pfactor} with all roots nonzero. For sufficiently
large $n$, it is positive. Write
\[
 S_j=\sum_{i=1}^s m_i(\ell_i/k_i)^j.
\]
Expanding the finitely many logarithms in
$P(n)=\rho n^d\prod_i(1+\ell_i/(k_i n))^{m_i}$ gives
\begin{equation}\label{eq:polylog}
 \log P(n)=\log\rho+d\log n
 +\sum_{j=1}^N\frac{(-1)^{j+1}S_j}{j n^j}
 +O(n^{-N-1}).
\end{equation}
Subtracting \eqref{eq:polylog} from \eqref{eq:logasym} proves an
all-order expansion for every integer sequence constructed here.
For example the first two correction terms are
\begin{equation}\label{eq:genericfirst}
 \frac{CR(n)}{P(n)}
 =\frac{CK_R}{\rho}\Lambda^n n^{-d-h/2}
 \left[1+\frac{\sigma_1-S_1}{n}
 +\frac{(\sigma_1-S_1)^2+S_2}{2n^2}+O(n^{-3})\right].
\end{equation}
For a negative leading coefficient, the same assertion holds with its
sign restored; the logarithmic derivation is applied to $|P(n)|$.

\subsection{The expansion for the certified sequence}
For $A$ one obtains
\[
 K_A=\frac1{\sqrt{60\pi}},\qquad \Lambda=\lambda,
 \qquad \sigma_1=\frac7{120}.
\]
In \eqref{eq:U}, $d=3$, $\rho=30$,
\[
 S_1=\frac12+\frac13+\frac15=\frac{31}{30},\qquad
 S_2=\frac14+\frac19+\frac1{25}=\frac{361}{900}.
\]
Therefore
\begin{equation}\label{eq:Uasym}
 \boxed{\displaystyle
 U_n=\frac{7}{5\sqrt{60\pi}}\,
       \frac{\lambda^n}{n^{7/2}}
       \left(1-\frac{39}{40n}+\frac{3893}{5760n^2}
       +O(n^{-3})\right).}
\end{equation}
All further coefficients are obtained by the explicit rational
operations in \eqref{eq:sigma}--\eqref{eq:polylog}. For $B$, the analogous
base constants are $\Lambda=2^{10}3^3$, $K_B=1/\sqrt{12\pi}$,
and $\sigma_1=1/36$. The method therefore gives comparably explicit
expansions for every companion sequence in the sharp-constant table.

\section{Verification record, limitations, and next questions}
\subsection{What was actually checked}
The archive contains two different kinds of computational work.
The finite certificates are part of the proof of the exact constants;
the regression tests merely provide additional protection against
implementation and transcription errors. The separate certificate
checker completed successfully for all 60 graphs, all 1,685 states,
and all 19,277 transition inequalities, including prime coverage and
sharpness witnesses.

The supplementary tests passed 90,298 comparisons between two forms
of the Landau step function, 808 direct rational checks of the eight
sharp quotients, and 480,048 small-prime valuation checks over
$0\leq n\leq10,000$. They also passed 3,336 direct checks of the
explicit uniform-product bounds, 4,000 valuation checks at reproducibly
sampled 256-bit indices, and 668 classification checks. The random
seed was 20260919. A further 203 exact checks verify the hypergeometric
recurrences and the first asymptotic coefficients. These tests are recorded in
\texttt{data/test\_results.json}; they are not cited as substitutes for
Theorem~\ref{thm:criterion}.

Every sporadic parameter list in the atlas was separately checked for
balance, height one, and nonnegativity at all rational breakpoints of
its Landau function. The atlas's parameter order follows Coserea's
OEIS table \cite{coserea}. The derived jump-sign columns are generated
from \eqref{eq:J}, not inferred from a finite list of sequence values.

\subsection{Limits of the conclusions}
The polynomial criterion applies to polynomials splitting over $\Q$.
It does not classify arbitrary irreducible quadratic or higher-degree
polynomials. Nor does the height-one sign rule extend unchanged to
higher height: there one must use the actual capacities in
\eqref{eq:capacity}, and multiplicities greater than one may occur.
The bound \eqref{eq:Cbound} is uniform and constructive, but often much
larger than the minimum; optimal constants for general products remain
a useful computational problem.

No Lean formalization was produced or run. The general proofs are
written in conventional mathematics, and the finite-graph certificates
are checked by ordinary exact-integer Python code. There is no reliance
on floating-point root finding, numerical asymptotics, or a conjectured
recurrence in the proof of integrality. A Lean development could split
naturally into the valuation identity, rational one-sided step lemmas,
the prime-power cutoff argument, and a generic finite-potential checker.

\subsection{Research directions arising from the proof}
A natural continuation is to classify nonsplit polynomial divisors.
Our necessity proof realizes prescribed rational approach points by
primes in arithmetic progressions. For higher-degree factors, roots
modulo primes and their distribution take the place of these simple
progressions, so the rational-root proof cannot simply be copied.

A second question is the behavior, as $r$ grows, of the \emph{optimal}
uniform multipliers for the two product families. The explicit lcm
bounds settle existence but do not predict the exact small-prime
exponents. The finite-state method suggests an effective route for
fixed $r$, with a separate challenge of controlling the state space
uniformly in $r$. Finally, the newly certified sequences invite
combinatorial models that would explain their integrality without
prime valuations; no such interpretation is assumed here.

\appendix
\section{Admissible denominators for all 52 sporadic ratios}\label{app:atlas}
In the following table, $a$ and $b$ are the numerator and denominator
parameter lists of $R(n)$. The column $K_-$ contains the $k$ with
$J(k)=-1$ and therefore describes primitive factors $kn+i$, $i>0$,
$(i,k)=1$. The column $K_+$ contains $J(k)=+1$ and describes $kn-i$.
The signs refer to jumps, not to the sign of the offset. Every allowed
root must still be simple. The parameter data are transcribed from
Coserea's table \cite{coserea}; the two $K$ columns are calculations
from Theorem~\ref{thm:heightone}. Row 1 is $B$ and row 31 is $A$.

\small
\begin{longtable}{rllll}
\toprule
Index & $a$ & $b$ & $K_-$ & $K_+$\\
\midrule
\endfirsthead
\toprule
Index & $a$ & $b$ & $K_-$ & $K_+$\\
\midrule
\endhead
\midrule\multicolumn{5}{r}{\small Continued on next page}\\
\endfoot
\bottomrule
\endlastfoot
\input{data/atlas_table.tex}
\end{longtable}
\normalsize

\Needspace{10\baselineskip}
\section{Reproducing the certificates and data}
The implementation requires Python 3.10 or later and uses only the
standard library. From the archive root, run
\Needspace{5\baselineskip}
\begin{lstlisting}[language=bash]
python code/verify_certificates.py
python code/test_research.py
\end{lstlisting}
The first command checks the supplied finite certificates without
calling the generator. To regenerate all finite graphs, equality
witnesses, sequence terms, and the atlas, run
\Needspace{5\baselineskip}
\begin{lstlisting}[language=bash]
python code/generate_artifacts.py
python code/verify_certificates.py
\end{lstlisting}
The checker deliberately raises an error if any state set is not
closed, any transition violates its potential inequality, a required
prime is missing, or a sharpness witness fails its direct valuation
calculation. The large-prime assertion is supplied by the proved
cutoff theorem and is not implemented as a search up to an arbitrary
prime bound.

The central library also exposes \texttt{Ratio}, \texttt{Factor},
\texttt{classify}, \texttt{one\_sided\_capacity}, and
\texttt{effective\_bound}. For example:
\Needspace{5\baselineskip}
\begin{lstlisting}[language=Python]
from factorial_divisibility import A, Factor, classify, effective_bound

assert classify(A, (Factor(2, 1), Factor(3, 1), Factor(5, 1)))
assert not classify(A, (Factor(4, 1),))
assert not classify(A, (Factor(2, 1, multiplicity=2),))
cutoff, valid_multiplier = effective_bound(A, (Factor(2, 1),))
# The returned multiplier is a proved bound, not necessarily minimal.
\end{lstlisting}
When using the library from the archive root, add \texttt{code/} to
\texttt{PYTHONPATH}, or run the example from within the code directory.
Nonprimitive factors must be normalized and their common integer
content treated separately. Equal roots must be merged by adding
multiplicities before calling the classifier. This avoids silently
mistaking two scalar multiples of a factor for two distinct roots.

The article compiles with ordinary pdfLaTeX and standard TeX Live
packages. Run \texttt{make pdf}, or run \texttt{pdflatex article.tex}
twice. The two atlas files and the certificate JSON are included, so
compilation does not require a network connection or regeneration.

\clearpage
\begin{thebibliography}{99}
\bibitem{oeisA}
The OEIS Foundation, \emph{A211417: $(30n)!n!/((15n)!(10n)!(6n)!)$}.
Comments by Peter Bala dated August 28, 2025; proof-status note by
Ralf Stephan dated June 30, 2026. Accessed September 19, 2026.
\url{https://oeis.org/A211417}.

\bibitem{oeisB}
The OEIS Foundation, \emph{A295431: $(12n)!n!/((6n)!(4n)!(3n)!)$}.
Uniform-divisibility conjectures by Peter Bala dated August 28, 2025.
Accessed September 19, 2026. \url{https://oeis.org/A295431}.

\bibitem{bala-note}
P. Bala, \emph{Proof of the integrality of A211417 and A211418}, April
2012; revised October 2015.
\url{https://oeis.org/A211417/a211417_1.txt}.

\bibitem{bober}
J. W. Bober, \emph{Factorial ratios, hypergeometric series, and a family
of step functions}, Journal of the London Mathematical Society (2)
\textbf{79} (2009), 422--444. Preprint arXiv:0709.1977.
\url{https://arxiv.org/abs/0709.1977}.
DOI: \href{https://doi.org/10.1112/jlms/jdn078}{10.1112/jlms/jdn078}.

\bibitem{adamczewski}
T. Adamczewski, \emph{OEIS Open: How many conjectures can language
models turn into theorems?}, arXiv:2608.11941, August 2026,
Appendix A, entry A211417.
\url{https://arxiv.org/abs/2608.11941}.

\bibitem{kedlaya}
K. S. Kedlaya, \emph{Notes on analytic number theory}, Chapter 4,
``Primes in arithmetic progressions,'' Theorem 4.2 and its proof.
\url{https://kskedlaya.org/ant/chap-primes-in-ap.html}.
Accessed September 19, 2026.

\bibitem{coserea}
G. Coserea, \emph{Table with the parameters of the 52 sporadic integral
factorial ratio sequences}, supplementary data for OEIS A295431.
\url{https://oeis.org/A295431/a295431.txt}.
Accessed September 19, 2026. Numerical parameter lists are reproduced
in the accompanying data with attribution; the admissibility columns
are computed in this work.

\bibitem{dlmf}
NIST Digital Library of Mathematical Functions, \emph{Section 5.11:
Asymptotic expansions}, especially formula 5.11.1 and the remainder
discussion. \url{https://dlmf.nist.gov/5.11}.
Accessed September 19, 2026.
\end{thebibliography}
\end{document}
